% Included by report/main.tex. This is an ordinary mathematical proof;
% the delivered PA194 endpoint is not claimed to be a newly compiled theorem.
\section{The PA-provable consistency equivalence}
\label{sec:pa194-calibration}

Let $M_{194}$ be the exact table \file{pa194/pa194.tm}, started in
state zero on the all-zero tape, and let $H_{194}$ assert that this
configuration has a finite halting trace. The halt label is $194$;
it is not one of the $194$ working states. Fix a conventional recursively
checked first-order PA proof predicate $\operatorname{Prf}_{\mathrm{PA}}$,
and put
\[
 \operatorname{Con}(\mathrm{PA})\equiv
 \neg\exists p\,\operatorname{Prf}_{\mathrm{PA}}
                 (p,\ulcorner0=1\urcorner).
\]

\begin{theorem}\label{thm:pa194-equivalence}
The delivered table satisfies
\[
 \mathrm{PA}\vdash
       \bigl(\neg H_{194}\leftrightarrow
                        \operatorname{Con}(\mathrm{PA})\bigr).
\]
This is an ordinary mathematical proof about PA-provability. The finite
certificates in the package verify its candidate-specific program facts;
the new endpoint has not been exported as a Lean or Rocq kernel theorem.
\end{theorem}

The change in logical source is the separate-diagonal emitter already
present in the supplied separate196 design. Its new association and
eight compatible physical folds give the present state count. We give
the changed command compiler explicitly, since the old nine-entry
compiler cannot be reused with its old indices. The commands below are
the original commands, including the test $t+h=1$; they are not the
later H-only Shift or MP-in-H variants.

\subsection{Exact numerical calculus and emitter}

Use Cantor pairing and the four constructors
\[
 p(a,b)=\frac{(a+b)(a+b+1)}2+b,\qquad I(A,B)=4p(A,B),
\]
\[
 M(a,b)=4p(a,b)+1,\quad U(a,B)=4p(B,a)+2,
 \quad E(a,b)=4p(a,b)+3.
\]
Thus $U$ deliberately transposes its arguments. Write
$\neg A=I(A,1)$, $A\wedge B=\neg(A\to\neg B)$, and
$\exists x_a A=\neg U(a,\neg A)$. Code $0$ is truth, while
$1=M(0,0)$ is the designated contradiction. In particular,
$I(0,0)=0$ is an actual numerical alias.

The six retained axiom families, in their original push order, are
\begin{align*}
 h_1(D,H)&=(\exists x_D\,U(D,H))\to H,\\
 h_2(D,A,H)&=U(D,U(A,H))\to U(A,U(D,H)),\\
 h_3(D,A,H)&=M(H,D)\to(E(H,A)\to M(A,D)),\\
 h_4(D,A,H)&=E(H,D)\to(M(A,H)\to M(A,D)),\\
 h_5(D,H)&=\exists x_D E(H,D),\\
 h_6(D,A,H)&=U(D,A)\to(U(D,A\to H)\to U(D,H)).
\end{align*}
Also retain the dominated propositional family
\[
 \rho_8(D,A,H,u)=((H\to u)\to D)\to
                        ((D\to H)\to(A\to H)),\qquad u\le D,
\]
modus ponens, and generalization $B\mapsto U(i,B)$ on closed
derivations. Here ``closed'' means without undischarged hypotheses;
the formula itself may contain free variable labels.

Put
\begin{align*}
 \Gamma(B)&=B\to U(B,B),\\
 V(a,b,c)&=M(a,b)\to U(c,M(a,b)),\\
 P(a,b,c)&=c\to M(c,b),\\
 S(a,b,c)&=U\bigl(b,(M(a,b)\to P(a,b,c))
                           \wedge(M(c,b)\to c)\bigr),\\
 Q(a,b,c)&=U(c,S(a,b,c)\to\neg P(a,b,c)),\\
 G(a,b,c)&=\neg(V(a,b,c)\to Q(a,b,c)).
\end{align*}
The separate calculus $\mathcal C_s$ adds every $\Gamma(B)$ and
every $G(a,b,c)$ with $c\ne a,b$. A failed guard emits $0$.
This adds no unlisted logical rule or arithmetic axiom. Truth is
derivable already: $\rho_8(0,0,0,0)=0$.

For an incoming store $(t,w,d,h)$, command $\mathsf{AX}(c,x,y)$ sets
$D=d+x$, $H=h+y$, $A=c$, and pushes, in order,
\[
 t,\ U(d,h),\ \rho_8(D,A,H,d),\ h_1,\ h_2,\ h_3,\ h_4,\ h_5,\ h_6,
 \ \Gamma(A).
\]
It leaves $G(H,D,A)$ in the top register if $A\ne H,D$, and
otherwise leaves $0$; both caches become zero. The list displayed
is push order, so its last entry is at the top of the resulting stack.
The complete ten-entry stack and final top are proved as identities
in a free term algebra, with all six incoming parameters symbolic,
by \file{mathematics/check_emitter_and_words.py}. Its grammar expansion
equals the archived separate196 word, and equals the PA195 word with
the \texttt{cons} at zero-based index $184$ deleted. This is a
190-token word. The final guard and cleanup blocks are unchanged.

\subsection{Proof equivalence with the calibrated diagonal calculus}

The PA195 calculus $\mathcal C_\Gamma$ has the same six heads,
$\rho_8$, and rules, but its guarded axiom is
\[
 F_\Gamma(a,b,c)=
       \neg\bigl(\Gamma(c)\to(V(a,b,c)\to Q(a,b,c))\bigr).
\]
The retained propositional basis and $h_1$ already derive the
classical implication and negation templates, independently of either
gate. In particular the restricted $\rho_8$ supplies the usual
positive basis and Peirce law, and $h_1$ supplies explosion.

\begin{lemma}\label{lem:pa194-calculi}
There are primitive-recursive, conclusion-preserving proof
translations in both directions between $\mathcal C_s$ and
$\mathcal C_\Gamma$, whose correctness PA proves.
\end{lemma}
\begin{proof}
Classical propositional templates extract $\Gamma(c)$ and
$\neg(V\to Q)=G$ from $F_\Gamma$. To supply an arbitrary
$\Gamma(B)$, use the guarded instance with $c=B$ and
$a=b=B+1$; its guard always succeeds. Conversely, proofs of
$\Gamma(c)$ and $\neg(V\to Q)$ give
$\neg(\Gamma(c)\to(V\to Q))$ by a fixed propositional template.
All other proof constructors are copied. These transformations recurse
on the finite derivation and use only code formation, concatenation,
and the displayed parameter choice. PA verifies the finitely many
templates and proves correctness by induction on that recursion.
No substitution into a freshness-restricted axiom is used.
\end{proof}

For completeness, we specify the logical calibration being inherited.
The preserved PA195 report (\file{reference/pa195_conpa}), in its
section \file{pa_calibration.tex}, constructs PA-verifiable proof
translations giving
\begin{equation}\label{eq:pa194-proofcalibration}
 \mathrm{PA}\vdash
 \bigl(\exists p\,\operatorname{Prf}_{\mathrm{PA}}
                  (p,\ulcorner0=1\urcorner)
       \leftrightarrow
       \exists d\,\operatorname{Prf}_{\Gamma}(d,1)\bigr).
\end{equation}
These are syntactic transformations, not an inference from the
existence of a standard model. In the forward direction the published
PA196 refutation compiler uses equality through direct and reverse
membership rewriting, neighborhood equality, and an equality witness.
The preserved diagonal proof replaces exactly the reverse and
neighborhood suppliers. Its negative-property specialization uses
direct rewriting, the retained equality witness, and membership
vacuity; it does not assume the equality-vacuity interface being
recovered. Choosing a fresh label as the code of the relevant formula
gives the required $\Gamma$ instance. Equality cases of numerical
labels are handled explicitly. These replacements are uniform finite
proof templates.

The reverse translation interprets $M(a,b)$ by the bit predicate
$\operatorname{Bit}(x_a,x_b)$, $E(a,b)$ by $x_a=x_b$, and $U$
by universal quantification. PA proves
$\operatorname{Bit}(s,z)\to z<s$ and hence
$\neg\operatorname{Bit}(z,z)$. Every label in a formula of code $B$
is strictly below $B$, so $\Gamma(B)$ is a valid vacuity instance.
The guarded separation supplier uses the PA proof template
\[
 \forall n\,\exists s<2^n\,\forall z<n\,
       (\operatorname{Bit}(s,z)\leftrightarrow\theta(z,\vec u))
\]
for an arbitrary formula $\theta$: induct on $n$, adding $2^n$
exactly in the case $\theta(n,\vec u)$. This procedure constructs
one induction instance syntactically; it does not evaluate $\theta$
or assume reflection for PA. When the source labels $a=b$, the
membership antecedent under that binder is
$\operatorname{Bit}(z,z)$ and is false; the preserved supplier uses
the corresponding singleton comprehension case. All substitutions
are capture avoiding. The truth alias is handled by the uniform
constructor equivalence
\[
 T(I(A,B))\leftrightarrow(T(A)\to T(B)),
\]
including $A=B=0$, rather than by an incorrect syntactic identity.
False code $1$ translates to $\operatorname{Bit}(x_0,x_0)$, whose
negation PA proves even under a binder for label zero.

Thus the retained head suppliers, generalization, and modus ponens
have uniform PA proof constructors. Induction on a coded derivation
checks the translation for every proof code. This proves the reverse
of (\ref{eq:pa194-proofcalibration}) internally, including nonstandard
codes in models of PA. The detailed constructions and their upstream
source references are preserved verbatim in the cited section.
Composing them with Lemma~\ref{lem:pa194-calculi} gives the same
biconditional with $\operatorname{Prf}_s(d,1)$.

\subsection{The exact ten-entry command compiler}

Let $[a]w=p(a,w)$ and let $\operatorname{hd},\operatorname{tl}$ be
Cantor projections. The unchanged commands act by
\[
 \mathsf S(t,w,d,h)=
 \begin{cases}
   \text{halt},&t+h=1,\\
   (h,\operatorname{tl}(w),d,\operatorname{hd}(w)),&t+h\ne1,
 \end{cases}
\]
and, for $v=d+x$,
\[
 \mathsf{MP}(x)(t,w,d,h)=
 \begin{cases}
   (v,[0]w,v,0),&I(t,v)=h,\\
   (0,w,v,0),&I(t,v)\ne h.
 \end{cases}
\]
The AX update is the ten-entry update above. Write $q_j$ for its
$j$th entry counted from the stack top. The reversal is explicit:
\[
\begin{array}{c|cccccccccc}
 j&1&2&3&4&5&6&7&8&9&10\\ \hline
 q_j&\Gamma(A)&h_6&h_5&h_4&h_3&h_2&h_1&\rho_8&U(d,h)&t.
\end{array}
\]
The gate and $q_1,\ldots,q_8$ have residue zero modulo four;
$q_9$ has residue two and is at least two. Directly applying Shift
therefore proves, without a restriction on the saved top $t$,
\begin{align}
 \mathsf{AX}(c,x,y)\mathsf S^9(t,w,d,h)
   &=(\rho_8(D,c,H,d),[t]w,0,U(d,h)),\label{eq:pa194-nine}\\
 \mathsf{AX}(c,x,y)\mathsf S^{10}(t,w,d,h)
   &=(U(d,h),w,0,t).\label{eq:pa194-ten}
\end{align}
Every tested pair before these endpoints has sum different from one.

Define
\begin{align*}
 \mathrm{pop}&=\mathsf{AX}(0,0,0)\mathsf S^{11},\\
 \mathrm{push}&=\mathsf{AX}(0,0,0)\mathsf S^9\mathsf{MP}(0),\\
 \mathrm{discard}&=\mathsf{MP}(1)\mathsf{AX}(1,0,0)
                                      \mathsf S^{10}\mathsf{MP}(0).
\end{align*}
Their exact identities are
\begin{align*}
 \mathrm{pop}(t,[a]w,d,h)&=(t,w,0,a),\\
 \mathrm{push}(t,w,d,h)&=(0,[t]w,0,0),\\
 \mathrm{discard}(t,w,0,0)&=(0,w,0,0).
\end{align*}
For pop, append one safe Shift to (\ref{eq:pa194-ten}). For push,
the MP comparison after (\ref{eq:pa194-nine}) compares an implication
code with a universal code, so fails by their different residues.
For discard, the first MP fails because $I(t,1)>0$, setting $d=1$.
Equation (\ref{eq:pa194-ten}) then gives $(U(1,0),w,0,0)$;
the final MP fails because $I(U(1,0),0)>0$.
These identities include $t=1$.

For $1\le j\le7$ define
\[
 \operatorname{select}_j(D,H,A)=
 \mathrm{discard}\,\mathsf{AX}(A,D,H)
                    \mathsf S^{j+1}\mathrm{pop}^{\,9-j}.
\]
It returns $(q_j,w,0,0)$ from $(t,w,0,0)$.
After the direct shifts its top is $q_j$, its cache is $q_{j+1}$,
and its remaining stack consists of $q_{j+2},\ldots,q_{10}$ and
the protected old $w$. Each pop preserves that top and consumes one
entry; the final cache is $q_{10}=0$ because discard saved truth.
Thus $\Gamma(B)$ uses $j=1,D=H=0,A=B$, while $h_1,\ldots,h_6$
use respectively $j=7,6,5,4,3,2$.

Let $\operatorname{offset}(i)$ be empty for $i=0$ and
$\mathsf{MP}(i)$ otherwise. From $(0,w,0,0)$ it sets $d=i$ and
changes nothing else. The dominated head is selected, for $u\le D$,
by
\[
 \mathrm{discard}\,\operatorname{offset}(u)\,
 \mathsf{AX}(A,D-u,H)\mathsf S^9\mathrm{pop}.
\]
The identity $u+(D-u)=D$ gives the desired $\rho_8$; the remaining
saved zero is removed by the final pop. The gate is selected by
$\mathrm{discard}\,\mathsf{AX}(c,b,a)\mathrm{pop}^{10}$.
It returns the guarded $G(a,b,c)$ and preserves $w$.

Generalization is implemented by
\[
 \mathrm{gen}_i=\mathrm{push}\,\operatorname{offset}(i)\,
                  \mathsf S\,\mathsf{AX}(i,0,0)\mathsf S^{10}.
\]
After push and offset, Shift transfers the saved formula $a$ to
$h$ without testing $a$; equation (\ref{eq:pa194-ten}) then returns
$(U(i,a),w,0,0)$. For modus ponens use
\[
 \mathrm{suffix}_b=
   \mathsf S\mathsf S\mathsf{MP}(b)\mathsf{AX}(b,0,0)\mathsf S^{11}.
\]
On $(0,[a][I(a,b)]w,0,0)$ it halts at the second Shift if $a=1$;
otherwise MP succeeds and yields $(b,[0]w,b,0)$, after which the
AX--Shift word returns $(b,w,0,0)$.

Compile a derivation by selecting its axiom, by applying
$\mathrm{gen}_i$ to a compiled generalization premise, and by
\[
 C(q)\,\mathrm{push}\,C(p)\,\mathrm{push}\,\mathrm{suffix}_b
\]
for premises $p:a$ and $q:I(a,b)$. Induction on the derivation
proves that $C(d)$, from every $(t,w,0,0)$, either halts or returns
$(a,w,0,0)$. A refutation therefore supplies the universally halting
word
\[
 W_d=\mathsf{AX}(0,0,0)\,C(d)\,\mathsf S.
\]
The first AX clears both caches from any incoming store; no reset
of the live stack or scheduler store is assumed. All identities are
equations of primitive-recursive arithmetic functions, with the
displayed residue and positivity proofs; PA proves their universal
closures and the compiler induction. The independent symbolic word
checker also replays these identities, including both MP-suffix cases.

\subsection{Scheduler, source binding, and exact physical traces}

For clarity, the existing enumeration needs no fairness assumption.
Encode a finite command word by continuation pointers:
\begin{align*}
 e([])&=0,\\
 e(\mathsf S::v)&=1+p(0,e(v)),\\
 e(\mathsf{MP}(x)::v)&=1+p(1,p(x,e(v))),\\
 e(\mathsf{AX}(c,x,y)::v)&=1+p(c+2,p(y,p(x,e(v)))).
\end{align*}
At a fetch boundary $(n,l)$ the scheduler decodes $n$ and increments
$n$ when $l=0$; otherwise it decodes $l-1$ and preserves $n$.
Its residual pointer is precisely $e(v)$ on the displayed words.
For any positive $l$, that residual is at most $l-1$, since each
Cantor projection is at most its input. Consequently each arbitrary
candidate finishes after a bounded number of commands or halts.
Fresh candidates enumerate $2,3,4,\ldots$ from startup.
Since $W_d$ begins with AX, $e(W_d)-1\ge3$, so its exact candidate
is eventually reached or an earlier command has already halted.
Bounded induction through the preceding candidates is a PA proof
of this assertion. Intermediate stores may be arbitrary: the
universal property of $W_d$ is exactly what is needed.

The source binding changes ordinal addresses explicitly:
\[
\begin{array}{l|rrr}
 &\text{upstream PA196}&\text{PA195}&\text{PA194}\\ \hline
 \text{halt site}&10&10&10\\
 \text{fetch entry}&25&25&25\\
 \text{AX entry}&50&50&50\\
 \text{MP entry}&3693&3581&3554\\
 \text{shared return}&3799&3687&3660\\
 \text{Shift entry}&3800&3688&3661\\
 \text{total primitive sites}&3812&3700&3673
\end{array}
\]
PA194 has $3504$ AX sites and the same $169$ exterior sites.
Deleting one \texttt{cons} macro removes $27$ primitive sites:
$11$ increments, $8$ decrements, and $8$ jumps. The independent
ordinal control-flow checks bind every successor, not only the
flattened token sequence, to the archived separate196 source.
Thus the arithmetic-macro composition and the exact command transfer
proof have the ten-entry AX effect above, with the displayed new
MP, return, and Shift addresses. The fetch convention and original
Shift/MP implementations remain identical under this ordinal map.

The final source program has $1457$ increment, $1107$ decrement, $1108$ jump,
and one halt primitive sites, giving
$1457+2\cdot1107+1108+1=4780$ symbolic branch contracts.
Both the unfolded and folded tables pass all of them. These quantify
over the entire register file and an arbitrary affine unary gap;
they are not $4780$ concrete numerical examples. The maximum source
path length is $38$, below header constant $48$.

Each symbolic contract consists of literal table rows, identical-tape
reexpressions, self-loop traversals of runs, and the inherited
register excursions. Induction on a run length or finite register
file proves each excursion's trace. The register-box rows are
preserved by the folds, and each accepted primitive trace has
strictly positive physical length and its exact stated endpoint.
Finite tapes, frame encodings, trace concatenation, and these
inductions are arithmetizable in PA. The candidate-specific finite
row equalities and completed symbolic traces therefore give a
PA-provable positive simulation uniformly in all coded frames.
The checked $58154$-step blank startup, including every nonzero
cell, ends at the clean fetch frame $(n,l)=(2,0)$ and store
$(0,0,0,0)$. Its fixed finite trace is also PA-provable.

\subsection{Sound stores and the two implications}

To reflect a possible physical halt, annotate both formula registers
$t,h$ and all canonical stack entries with PA proofs of their
translations $T$. The stack requires care: $p(0,0)=0$, so it is not
a free list encoding. For $w>0$, $\operatorname{tl}(w)<w$; define
$\delta(w)$ as the least $k\le w$ for which
$\operatorname{tl}^{k}(w)=0$. The annotation has exactly
$\delta(w)$ entries, certifying the successive heads. A push
prepends a proof except when both the code and stack are zero;
then the canonical stack remains empty. Popping zero returns zero
with its truth proof. The identities and compiler above also cover
these cases, including the alias $I(0,0)=0$.

AX preserves this invariant by the PA suppliers already specified,
including its unconditional separate $\Gamma$ even on a failed
guard. A successful MP combines the proofs of $t$ and $h=I(t,v)$;
its pushed zero receives a truth proof. A failed MP returns truth.
A nonhalting Shift transfers the corresponding annotations.
If Shift halts, $t+h=1$ implies $t=1$ or $h=1$, so one annotation
and the fixed PA proof of $\neg T(1)$ produce a PA refutation.
Each update is primitive recursive and PA verifies its invariant.

Work now in PA under $\operatorname{Con}(\mathrm{PA})$.
No annotated command can halt. The exact source command contract
therefore supplies a positive execution to the next annotated fetch
boundary. Given a proposed physical halting time $T$, primitive
recursion through $T+1$ command boundaries supplies more than $T$
primitive source steps, and concatenating the positive primitive
contracts supplies a physical execution longer than $T$.
All these are bounded coded iterations with explicit finite
witnesses. Determinism contradicts the proposed halt: a partial run
cannot pass through the halt label and later reach a nonhalt frame.
An alleged halt inside startup contradicts its fixed checked trace.
Thus PA proves $H_{194}\to\neg\operatorname{Con}(\mathrm{PA})$.
This argument needs neither inverse decoding of an arbitrary tape
nor a least-time search for a future frame.

Conversely, a PA refutation gives a $\mathcal C_s$ refutation by
(\ref{eq:pa194-proofcalibration}) and Lemma~\ref{lem:pa194-calculi}.
Its compiled $W_d$ halts from every store. The bounded scheduler
argument reaches that word or halts earlier; the arithmetic-macro
and positive physical trace contracts then give a finite blank-tape
halt of the literal table. All proof-producing operations and
correctness inductions just used are PA-verifiable, so
$\mathrm{PA}\vdash\neg\operatorname{Con}(\mathrm{PA})\to H_{194}$.
The two implications prove Theorem~\ref{thm:pa194-equivalence}.
